% commentaire.tex — f169v, vv. 590–634 (conversion initiale des travaux antérieurs)
\comm{593}{ἠλίθιος ὅστις τιλλόμενος ἠνείχετο;}{Willems, chevalier servant du \italique{Ravennas}, défend encore une fois la leçon du manuscrit dans sa \nieme{18} note, à retrouver à la fin.}
\comm{602}{ὦ πρόξενε.}{Les proxènes ou hôtes publics, agents officiels dont les fonctions consistaient principalement à recevoir et à héberger les délégués de l'État dont ils exerçaient la proxénie. (Note de Willems.)}
\comm{605}{Κλεωνύμου γυνή.}{L'évocation de Cléonyme est l'occasion de dresse un bref résumé des railleries à l'encontre de ce personnage parmi les plus plaisantés de l'Ancienne comédie. Cet Athénien nous est autrement connu par trois décrets qu'il proposa à la \grec{βουλή}\nc{\italique{Inscriptiones Graecae}, I\expo{3}, 61.34, 68.5, 69.3-4.} l'an 426/5 ainsi qu'une proposition de fixer à 1000 drachmes la récompense contre des informations touchant à l'affaire de la parodie des Mystères d'Éleusis\nc{Andocide, I-\italique{\grec{Περὶ τῶν μυστηρίων}}, 27. La démagogie de Cléonyme se voit elle-même débordée par celle de Pisandre qui renchérit à 10000 drachmes !..} (contemporaine de celle des Hermocopides, l'an 415). Cléonyme est victime d'un tir de barrage d'Aristophane portant à la fois :\\
— au moral : sa lâcheté (il lui était arrivé un jour de jeter son bouclier pour mieux fuir. C'est surtout là-dessus que le poète ne tarit point : \italique{Cavaliers} v. 1372, \italique{Nuées} v. 353, \italique{Guêpes} v. 19, 592, \italique{Paix} v. 446, 672, 1186 (non nommé), 1275, 1295, \italique{Oiseaux} v. 290, 1475, 1481) ; sa duplicité (\italique{Acharniens} v. 88, il est traité de parjure v. 399 des \italique{Nuées}, de flatteur au v. 592 des \italique{Guêpes} (il est appelé Κολακώνυμος), de sycophante au v. 1479 des \italique{Oiseaux}) ; ses manières brusques de grand dadais (\italique{Acharniens} v. 844) ; ses mœurs peu ragoûtantes (accusé d'être un \grec{καταπύγων} au v. 675 des \italique{Nuées}) ;\\
— au physique : une laideur repoussante (v. 822 des \italique{Guêpes}, ) ; son androgynie (Socrate, dans une intention de parler exactement, emploie le féminin \grec{Κλεωνύμη} au v. 680 des \italique{Nuées}) ; une taille excessive (\italique{Acharniens} v. 88, \italique{Oiseaux}, v. 1476) ; sa gloutonnerie (\italique{Cavaliers} v. 1290, \italique{Oiseaux} v. 289).\\
Jouant sur les contrastes, Aristophane nous présente une femme de Cléonyme \italique{petite} de nom (c'est le sens de \grec{Μίκα}), virile et hardie dans les débats.\\
Revenons sur le grief le plus reproché à Cléonyme : avoir jeté son bouclier pour mieux prendre ses jambes à son cou, sa \grec{ῥίψασπις}. Il n'y a rien à y redire : le crime était théoriquement puni de mort\nc{Andocide, I-\italique{\grec{Περὶ τῶν μυστηρίων}}, 74.} et l'on comprend aisément que le poète ait exposé à l'opprobre publique le couard qui avait encore l'outrecuidance de venir palabrer après son forfait. Exégèse par trop simple pour nos éminents professeurs – dont le délassement consiste à extravaguer sur les passages les plus obscurs et remettre systématiquement en doute les passages les plus clairs\nc{Willems, tome I, p. 108, s'en plaignait déjà en son temps.}. On ne la fait pas à ces \italique{viri docti}, comme disait Fritzsche : le dernier en date\nc{I. C. Storey, «~The `Blameless Shield' of Kleonymus~», in \italique{Rheinisches Museum für Philologie}, n°132 (1989), p. 247-261.} argue de ce que Cléonyme serait mort – condamné à mort – bien avant les dernière traces archéologiques de sa vie démocratique s'il avait effectivement commis ce qu'Aristophane lui reproche... avec le même raisonnement, les historiens des prochains millénaires déduiront que l'\italique{Élysée}, \italique{Matignon} et le \italique{palais Bourbon} ne pouvaient désigner que des prisons au vu des canailles hors-la-loi qui les fréquentaient ; prisons flanquées, semblera-t-il, d'une maison de santé pour ceux dont la panse menacait de crever à tout instant : le \italique{palais du Luxembourg}\nc{``Palais'' est à prendre dans son acception judiciaire ; suite logique des \italique{palais} de justice.}. Mais il n'y a pas besoin d'extrapoler : le même raisonnement intenable entraînait déjà de refuser le vers 1166 des \italique{Acharniens}\nc{voir Willems, t. I p. 106 sq.} : comment un bandit (un \italique{Oreste}, nom générique ; un \italique{Cartouche} ou \italique{Mandrin} chez nous) aurait-il pu sévir la nuit à Athènes alors que le banditisme était aussi puni de mort ?..
Certes, « \anglais{in interpreting the personal jokes in Aristophanes one must always be wary of the too serious or too literal explanation of the text} »\nc{«~The `Blameless Shield' of Kleonymus~», p. 248.}, mais l'\italique{on ne détruit que ce que l'on remplace}, et pour remplacer une explication qui coule de source comme celle de nos passages, il vaut mieux présenter de solides arguments. Or M. Storey ne me fera pas accroire sans preuve\nc{Car je tiens pour nulle et non avenue sa citation du passage des \italique{Cavaliers} (v. 1369-1372) qui \italique{prouve} simplement que M. Storey n'a pas saisi l'image et le jeu de mots d'Aristophane sur \grec{κατὰ σπουδὰς} (mis à dessein pour faire entendre κατ᾿ ἀσπίδας en lieu et place de l'expression figée `κατὰ σπουδήν' qu'on trouve chez Thucydide, Xénophon, Élien, Plutarque, etc.). Au reste, si M. Storey convient que les vers 827-829 des \italique{Thesmophories} visent Cléonyme, je lui saurais gré d'expliquer le tour `ἀπὸ τῶν ὤμων' à l'aune de sa brillante exégèse.} que le motif pour lequel Aristophane poursuivit inlassablement de sa verve pendant des années Cléonyme fût, non pas une désertion devant l'ennemi doublée d'un abandon d'armes, mais une banale brigue pour échapper au service militaire\nc{Si banale qu'on en trouve à la pelle parmi nos modernes brigands. Citons seulement l'exemple le plus caricatural : M. Ciotti-le-bien-nommé, alors qu'il n'était encore que vil attaché-parlementaire fit jouer ses relations pour échapper au service militaire... maintenant que son derrière est confortablement installé au palais bourbeux, bien loin des fangeuses casernes, Chiotti n'a de cesse de haranguer le veau français pour un rétablissement dudit service militaire.}.}
\comm{609}{τίτθη.}{Remarquons en passant qu'une nourrice à Athènes n'était pas nécessairement une esclave. Cf. Démosth., LVII, 35. (Note de Willems.)}
\comm{616}{κάρδαμον}{Les botanistes s'accordent, je crois, à reconnaître dans le κάρδαμον, non le \italique{cresson}, comme on l'a cru jusqu'en ces derniers temps, mais une autre crucifère, connue des savants sous le nom de \latin{erucaria aleppica}. Suivant Fraas, \italique{Synopsis florae classicae}, les Grecs lui donnent encore aujourd'hui le même nom. Hippocrate dit positivement que le \grec{κάρδαμον} provoque la dysurie (\italique{Du Régime}, II, 54). On cite à ce sujet une singulière anecdote, recueillie par Xénophon. Chez les Perses, où il passait pour indécent de cracher, se moucher ou s'écarter pour faire de l'eau, on ne donnait aux enfants à leur repas du midi que du pain, avec du cardame pour tout assaisonnement. (\italique{Cyrop.}, I, 2, 8, 11 et 16) Mettez cela d'accord avec les traités de thérapeutique où le cresson est donné pour un sialagogue. (Note de Willems. Les dictionnaires  \LSJ , Bailly et Pape désignent le \grec{κάρδαμον} respectivement comme « \anglais{nose-smart}, \latin{Lepidium sativum} »,\italique{cresson alénois} ; « cresson » et « \allemand{eine Art Kresse mit bitterem Kraut} ».Pour aller dans le sens de Willems, ajoutons que le cresson, outre ses vertus sialagogiques, est loin d'être dysurique : il a des vertus au contraire diurétiques. Cf. le \italique{Dictionnaire des plantes médicinales et vénéneuses de France} de Paul-Victor Fournier, p. 731. Mais la raison de sa consommation chez les Perses est peut-être à chercher ailleurs : cette plante était réputée donner du courage.)}
\comm{619}{τὸν ἐκ Κοθωκιδῶν;}{« Du $[$dème$]$ des Cothôcides ». Ce dème faisait partie de la tribu des \grec{Οἰνηίς} au nord-ouest de la ville\nc{L'on pourra consulter le travail de John S. Traill, \italique{The Political Organization of Attica}, \italique{Hesperia} supplement XIV (1975). Voir notamment les pages 49, 62 et 68 ainsi que les cartes d'Athènes à la fin.}. Y a-t-il une raison particulière pour que Mnésiloque évoque ce dème ? certainement... mais cette raison nous reste obscure et aucune explication satisfaisante n'a jusqu'ici été avancée.}
\comm{624}{συσκηνήτρια.}{Il faut croire d'après ce passage et le vers 658 que lors des Thesmophories, des tentes pour deux femmes étaient installées à proximité du Pnyx.}
\comm{633}{σκάφιον Ξένυλλ᾽ ᾔτησεν· οὐ γὰρ ἦν ἀμίς.}{Acculé, c'est à ce vers que Mnésiloque fait son hara-kiri : il lance une dernière boutade en guise de baroud d'honneur ! L'explication communément admise laisse songeur quant à l'humour des commentateurs d'Aristophane : Mnésiloque emploierait selon eux un terme impropre (l'\italique{amis} serait une sorte de pot de chambre à usage exclusivement masculin en raison d'un col étroit... mais sa forme varie et le terme pourrait très bien aussi désigner ce genre d'ustensile pour les femmes\nc{Pour toutes les remarques de cette note, je renvoie aux articles \latin{amis} et \latin{scaphium} du \italique{Dictionnaire des antiquités grecques et romaines} de Daremberg \& Saglio.}). Quoi qu'il en soit, c'est un peu \italique{gros} : tout rustre qu'il soit, Mnésiloque ne peut pas commettre une telle bévue ; je dirais même plus~: dans ce domaine, il est certainement tout aussi savant que l'\italique{Athénienne de la rue}. 
Voici une explication bien meilleure : le \grec{σκάφιον} désignait à l'origine une \italique{coupe} (en forme de \italique{barque} selon l'étymologie). Il n'acquiert le sens quasi-exclusif d'\italique{urinoir} qu'en passant au latin \latin{scaphium}. À l'époque d'Aristophane, les deux sens devaient encore co-exister. La première partie du vers « σκάφιον Ξένυλλ᾽ ᾔτησεν » était donc interprétée par Mica favorablement~: Xénylla a demandé une coupe $[$pour faire une libation$]$... la chute est rude et Mnésiloque dans la seconde partie accuse les femmes d'un usage détourné des coupes consacrées qui passe la mesure : cet individu n'a rien des manières d'une femme : après avoir uriné publiquement, voilà qu'il prononce maintenant des \italique{insanités} (οὐδὲν ὑγιὲς v. 636).\\
Voici pour le sens propre -- si l'on peut dire. Mais, comme de coutume, Aristophane joue à différents niveaux : la réplique est aussi un clin d'œil au proverbe grec bien connu\nc{Il a été relevé par Suidas en α 1504 et α 1574. Érasme y consacre une entrée dans ses \italique{Adages} (au n°1149) et l'évoque aussi au n°1880. Maintes fois recensé par les parémiographes (par exemple Zénobius I, 83), il était connu de Plutarque (\italique{Du bavardage}, OM, 6, 35, 512e) et le scholiaste de la \italique{Paix} d'Aristophane le cite encore à propos de la chanson de Datis le Lydien aux vers 289-291 et de l'équivoque en grec que suggère la forme barbaristique χαίρομαι -- au passage : Aristophane joue sur le stéréotype du débauché « Lydien ($=$ Datis) à midi », Λυδὸς ἐν μεσημβρίᾳ (rapporté par Athénée, XII, 515d-517a et recensé par Érasme dans ses \italique{Adages}, n° 1594). Enfin, le scholiaste de Théocrite (\italique{Idylles}, IV, 10) cite aussi de ce proverbe. Pour finir, Aristophane y fait de nouveau allusion au fr. 927 Kassel-Austin : ἄγροικός εἰμι· τὴν σκάφην σκάφην λέγω.} :\citComm{Ἄμας ἀπῄτουν, οἱ δ' ἀπηρνοῦντο σκάφας\\
Je demandais des bêches, ils me refusent des seaux.}(Jeux de consonances entre ἄμη/ἀμίς et σκάφη/σκάφιον.)
Apprécions à sa juste valeur l'à-propos de l'allusion au proverbe : il repose sur un quiproquo dû à la polysémie d'ἄμη -- comme est polysémique le σκάφιον dont parle Mnésiloque. De plus, le proverbe évoque ceux qui fournissent des réponses sans rapport avec la question posée -- ce qui est aussi le cas de notre pauvre Mnésiloque, bien embarrassé.
ἄμη et σκάφη pouvaient en effet être tous deux synonymes de \italique{seau}\nc{Comme dans le proverbe pour évoquer abondance de biens : ἄμαις καὶ σκάφαις ἀρύσασθαι. Abondance de biens nuit, soit dit en passant, à la santé du corps comme celle de l'esprit. Bloy n'a, à ma connaissance, étrangement pas fait l'exégèse de ce lieu commun -- Flaubert aurait parlé d'une idée reçue. Les latins et leur solide bon sens savaient que \latin{Copia parit fastidium} et \latin{Neque laus in copia, neque culpa in penuria constitit}.)}... mais c'était des \italique{houes} (autre sens de ἄμη) que demandait le paysan du proverbe...}
